\documentclass[12pt]{article}
\usepackage[utf8]{inputenc}
\usepackage[T1]{fontenc}
\usepackage{hyperref}
\title{Southern Pame Language Resources for K-12 Curriculum Development}
\author{Compiled by Claude Sonnet 5.5 (claude-sonnet-5-5) for Andr\'as Kornai}
\date{}
\begin{document}
\maketitle
\section{Introduction}
Southern Pame (ISO 639-3 pmz) is described as extinct; its speakers were in Hidalgo, Mexico. The name ``Southern Pame'' is used inconsistently in sources, and searches found no material specific to this language. The two links below concern the Pame people and the related Northern Pame, not Southern Pame itself.
\section{Primary Educational Resources}
None found.
\section{Linguistic and Reference Materials}
\begin{itemize}
\item Pame people, overview of the Pame groups and their names: \url{https://en.wikipedia.org/wiki/Pame_people}
\item Northern Pame (a different, living language), Endangered Languages Project sample page: \url{https://www.endangeredlanguages.com/lang/4978/samples/6401}
\end{itemize}
\section{Religious and Community Resources}
None found.
\section{Digital and Audio Resources}
None found for Southern Pame.
\section{Curriculum Development Context}
With no documented speakers or materials found, a heritage curriculum would have to draw on archival work and on the better documented Northern and Central Pame, with the community's involvement.
\end{document}